\chapter{Discussion}
\label{ch:discussion}

\section{What the results jointly show}

The results describe a judge ecosystem with useful but bounded reliability. The RQ1 agreement level shows that automated judgments capture a meaningful portion of the human preference reference pattern, yet leave a substantial disagreement rate. RQ2 shows that the same frozen prompt/protocol is usually repeated consistently. RQ3 demonstrates why this is not enough: changing only answer order can change a decisive mapped outcome for a non-trivial share of eligible pairs. This combination is precisely why a single global ``reliability'' label would be misleading.

The RQ4 and RQ5 null results are also informative when interpreted with the right scope. The specified redundant-text and presentation/list-prefix treatments did not produce stable wins in their eligible controlled pairs. That finding limits a claim about these treatments; it does not settle a broader debate about verbosity, style, salience, or prompt-dependent formatting. A scientifically useful null result is not an invitation to say ``no bias exists.'' It is a well-bounded observation under a frozen intervention.

RQ6 adds another dimension. The aggregate same-family preference rate is close to 50\%, and its interval spans 50\%, so a uniform overall preference claim is not supported. At the same time, judge-level patterns differ sharply. This is a caution against interpreting aggregate evidence as an explanation of every component judge. It also illustrates why counterbalancing, exact answer identity, and reporting coverage are essential in source-family questions.

\section{Interpreting mitigation}

DUAL\_SWAP and Multi-Judge Consensus address different practical concerns. DUAL\_SWAP asks whether a within-judge decision remains compatible with a presentation-consistency requirement. It may remove decisions whose evidence is less stable under answer swapping; in doing so it visibly reduces coverage. In the final analysis its matched agreement change is small and its confidence interval crosses zero. This does not make the method useless. It means the particular agreement-gain estimand is uncertain in this population and should always be read with the 21.03 point coverage reduction.

Multi-Judge Consensus aggregates four independent judge votes under a strict frozen protocol. On its own retained pairs, it improves agreement with human preference reference labels by 5.02 points compared with an equal-weight individual-judge baseline. The confidence interval is positive under the frozen analysis. This is evidence for a secondary mitigation result, not evidence that an ensemble is universally better. The protocol excludes pairs without the required valid votes, retains 69.83\% of its planned population, and evaluates a cross-judge aggregation family rather than a within-judge filter.

The two results should not be placed on a one-dimensional race track. Their comparators differ. Their retained units differ. Their operational costs and failure modes differ. A future applied deployment may choose a DUAL\_SWAP-like filter when the priority is presentation robustness within a judge, or a multi-judge procedure when the priority is aggregation across independent models and the coverage/cost profile is acceptable. That is a decision framework, not a claim of universal superiority.

\section{Threats to validity}

\subsection{Construct validity}
Human preference reference labels are a defined reference target, not ground truth. They may reflect preference, helpfulness, style, or task conventions. Exact agreement and $\kappa$ quantify correspondence with that target. They do not establish factual correctness, safety, or utility in every deployment. Similarly, repeated-evaluation consistency measures repeatability under a fixed condition, not an abstract stochastic property across arbitrary parameter settings.

\subsection{Internal validity}
The source-text remediation strengthens internal validity by binding controlled records to exact source text and preserving a documented source-precedence policy. The study avoids imputation, model substitution, and parser relaxation. Nevertheless, terminally unavailable observations remain. Fourteen terminal Claude observations were not treated as missing completely at random or missing at random. Complete-case estimators preserve transparency but may differ from an unavailable full outcome population.

\subsection{External validity}
The population consists of 80 MT-Bench questions represented as two turns, exact candidate answers, four judge families, and frozen protocols. Results should not automatically transfer to other tasks, languages, model releases, prompt templates, scoring rubrics, or provider routing configurations. The narrow RQ4/RQ5 treatments are especially unsuitable for generalized claims about all lengthy or formatted responses.

\subsection{Statistical conclusion validity}
The report gives denominators, coverage, and bootstrap intervals to prevent over-reading a point estimate. The DUAL\_SWAP interval crossing zero is important. The RQ6 aggregate interval spanning 50\% is important. Multi-Judge's positive interval is important, but it is conditional on a strict retained population and its own comparator. Multiple related descriptive breakdowns are presented as descriptive rather than as a series of independent confirmatory discoveries.

\subsection{Operational validity}
Provider failures, invalid responses, and persistence failures can distort an experiment if they are silently retried, relabeled, or dropped. The project records durable states and distinguishes physical completion, valid outcomes, and analytical eligibility. Live manual evaluation has a separate server-side safety gate and a separate telemetry class. A successful live manual trial cannot modify the frozen controlled evidence.

\section{Practical recommendations}

First, do not report an LLM evaluator only through one agreement percentage. Pair agreement with repeatability, order sensitivity, missingness, and coverage. Second, retain exact answer provenance; source labels are not a substitute for source identity. Third, expose the denominator and the rule that created it. Fourth, prefer protocols whose final API can fail closed rather than return a convenient hard-coded fallback. Fifth, distinguish a mitigation's comparator and retained population before comparing its effect to another method. Finally, preserve evidence packages and release snapshots so that a final dashboard value can be independently recomputed later.

\section{Future work}

Future work could use a larger and more diverse prompt population, independently refreshed human preference labels, preregistered primary/secondary analyses, alternative consensus rules, and cost-aware operational comparisons. It could also study the source of the RQ3 pattern through controlled judge prompts and model versions. Such work should preserve the present project's discipline: exact source text, explicit mapping of swapped presentations, visible terminal states, and modest language when causal explanation is not identified.
